\clearpage
\VerdierPhysicalPageAnchor{0245}%
\begingroup\small\setlength{\parskip}{0pt}\setlength{\emergencystretch}{1.5em}
\setlength{\abovedisplayskip}{4pt plus 2pt minus 2pt}%
\setlength{\belowdisplayskip}{4pt plus 2pt minus 2pt}%

\VerdierSemanticAnchor{ch4d-4-6-3-b}
\noindent\textit{b) Let
$Q\colon K(\mathcal A)\longrightarrow D(\mathcal A)$ be the quotient
functor. The complex $I^\bullet$ is $Q$-free on the right (Chapter~II,
2.3.3); i.e.\ for every complex $Y$, the canonical morphism}
\[
\operatorname{Hom}_{K(\mathcal A)}(Y,I^\bullet)
\longrightarrow\operatorname{Hom}_{D(\mathcal A)}(Y,I^\bullet)
\]
\textit{is an isomorphism.}

\VerdierSemanticAnchor{ch4d-4-6-3-c}
\noindent\textit{c) Let}
\[
H^*(I^\bullet)\colon\cdots\longrightarrow H^p(I^\bullet)
\longrightarrow H^{p+1}(I^\bullet)\longrightarrow\cdots
\]
\textit{be the complex whose components are the $H^p(I^\bullet)$,
$p\in\mathbb Z$, and whose differentials are zero. There is a unique
isomorphism in $K(\mathcal A)$
$I^\bullet\longrightarrow H^*(I^\bullet)$ inducing the identity on the
cohomology objects.}

\VerdierSemanticAnchor{ch4d-4-6-3-d}
\noindent\textit{d) The canonical morphisms
$H^*(I^\bullet)\longrightarrow H^p(I^\bullet)[-p]$ define isomorphisms in
$K(\mathcal A)$ and in $D(\mathcal A)$:}
\[
H^*(I^\bullet)\xrightarrow{\sim}
\prod_{p\in\mathbb Z}H^p(I^\bullet)[-p].
\]

\VerdierSemanticAnchor{ch4d-4-6-3-proof-abcd}
We show that a) implies b). Indeed, it suffices (Chapter~II, 2.3.3) to show
that for every acyclic complex $Y$,
$\operatorname{Hom}_{K(\mathcal A)}(Y,I^\bullet)=0$, which follows from a).
We show that a) implies c). Indeed, it suffices to replace, in the
isomorphism of a), the complex $Y$ by $I^\bullet$ and the complex
$I^\bullet$ by $H^*(I^\bullet)$. One sees likewise that a) implies d) by
replacing, in the isomorphism of a), the complex $I^\bullet$ by the complex
$H^*(I^\bullet)$. We prove a). It follows immediately from the definitions
that the complex $I^\bullet$ is isomorphic to a complex of the form

\VerdierSemanticAnchor{ch4d-4-6-3-proof-a-decomposition}
\[
\resizebox{0.98\linewidth}{!}{$
\displaystyle
\cdots\longrightarrow
J^{n-2}\oplus H^{n-1}\oplus J^{n-1}
\xrightarrow{\left(\begin{smallmatrix}
0&0&\mathrm{id}\\0&0&0\\0&0&0
\end{smallmatrix}\right)}
J^{n-1}\oplus H^n\oplus J^n
\xrightarrow{\left(\begin{smallmatrix}
0&0&\mathrm{id}\\0&0&0\\0&0&0
\end{smallmatrix}\right)}
J^n\oplus H^{n+1}\oplus J^{n+1}
\longrightarrow\cdots
$}.
\]

\VerdierSemanticAnchor{ch4d-4-6-3-proof-a-homotopy}
Consequently, there is a homotopy equivalence between $I^\bullet$ and the
complex $H^*(I^\bullet)$. To prove a), we may therefore suppose that
$I^\bullet=H^*(I^\bullet)$. It is then clear that there is an isomorphism

\VerdierSemanticAnchor{ch4d-4-6-3-proof-a-hom}
\[
\operatorname{Hom}_{K(\mathcal A)}\bigl(Y,H^*(I^\bullet)\bigr)
\xrightarrow{\sim}
\prod_{p\in\mathbb Z}
\operatorname{Hom}_{K(\mathcal A)}
\bigl(Y,H^p(I^\bullet)[-p]\bigr).
\]
Moreover, since the objects $H^p(I^\bullet)$ are injective, the canonical
morphism
\[
\operatorname{Hom}_{K(\mathcal A)}
\bigl(Y,H^p(I^\bullet)[-p]\bigr)
\longrightarrow
\operatorname{Hom}_{\mathcal A}\bigl(H^p(Y),H^p(I^\bullet)\bigr)
\]
is an isomorphism.
\endgroup
